\documentclass[10pt,a4paper]{amsart}
\usepackage[margin=19mm]{geometry}
\usepackage{amsmath,amssymb,mathtools}
\usepackage[scr=rsfs,cal=euler]{mathalfa}
\usepackage{xfrac}
\usepackage[unicode,hidelinks,bookmarks=false]{hyperref}
\newcommand{\ie}{i.e.\ }
\newcommand{\cf}{cf.\ }
\newcommand{\resp}{respectively\ }
\newcommand{\Subsection}{\subsection}
\providecommand{\nobreakdash}{\nobreak\hbox{-}}
\setlength{\emergencystretch}{2em}
\multlinegap0pt
\usepackage{bm}
\usepackage{bbm}
\usepackage{dsfont}
\usepackage{xspace}
\usepackage{cancel}
\usepackage{mathtools}
\usepackage{cleveref}
\usepackage{amsthm}
\makeatletter
\@ifundefined{thm}{\newtheorem{thm}{Thm}}{}
\@ifundefined{lem}{\newtheorem{lem}[thm]{Lemma}}{}
\makeatother
\def\E{\mathbb{E}}
\def\N{\mathbb{N}}
\def\P{\mathbb{P}}
\def\R{\mathbb{R}}
\title[Quenched law of large numbers and quenched central limit the]{Quenched law of large numbers and quenched central limit theorem for multi-player leagues with ergodic strengths}
\author{Jacopo Borga, Benedetta Cavalli}
\date{Source: arXiv:2103.01320v1 (CC BY 4.0)}
\begin{document}
\maketitle

\noindent\emph{Source and licence.} Quenched law of large numbers and quenched central limit theorem for multi-player leagues with ergodic strengths. Version arXiv:2103.01320v1 (CC BY 4.0), distributed under the Creative Commons Attribution 4.0 International licence, as recorded in the arXiv licence metadata for this exact version and shown on the abstract page https://arxiv.org/abs/2103.01320v1. Full rights notice: licenses/arxiv-2103.01320-CC-BY-4.0.txt.\par\smallskip

\noindent\textit{Editorial scope.} Counted unit: the Lem whose source label is \texttt{lem:ind\_fnct} in the article (source0.tex lines 723--728, source label \texttt{lem:ind\_fnct}), delivered together with the article's own complete original proof of it (source0.tex lines 730--818, one \texttt{proof} environment of 89 lines). No mathematics is written, repaired or re-derived: statement and proof are the article's own text line for line. Required context retained uncounted: lem:first\_sec\_mom\_for\_ret (lines 656--661). Every definition, display and label of those lines is retained unchanged. Omitted from this package and not redistributed: 1--655, 662--722, 729--729, 819--1296. Nothing inside a retained excerpt is omitted or altered: every retained body is delivered exactly as the article's lines are written, so no omission receipt of any kind accompanies this package. Citations: the retained excerpts carry no citation command at all, so no bibliography entry of the article is transcribed and no reference is redistributed. Reference adaptation: none was needed and none was made; every reference inside the delivered unit resolves inside it, so no unresolved reference or citation is produced. Printed slips kept literal: no printed word, symbol, hypothesis or number of the counted statement or of its proof was altered. Layout and class: the article's own class is \texttt{article} with packages including geometry, xcolor, hyperref, babel, authblk, amsmath, graphicx, amsthm, cleveref, autonum, bm, amsfonts, mathtools, dsfont; the wrapper is the lane's pinned \texttt{amsart} editorial wrapper at 10pt on A4, which restates the theorem-like environments the retained text needs and adds the article's own macro definitions that the retained text needs (\texttt{\textbackslash E}, \texttt{\textbackslash N}, \texttt{\textbackslash P}, \texttt{\textbackslash R}), with extra wrapper packages (bm, bbm, dsfont, xspace, cancel, mathtools, cleveref). The delivered numbering is the unit's own. Assets: no retained span contains a figure environment, a graphics-inclusion command, a PDF-page inclusion, a table or a TikZ picture; no image file is copied into this package, the package declares no asset dependency and no image file is redistributed with it. Licence: version arXiv:2103.01320v1 of this article is distributed under the Creative Commons Attribution 4.0 International licence, as recorded in the arXiv licence metadata for this exact version and shown on the abstract page https://arxiv.org/abs/2103.01320v1. A full rights notice is delivered at licenses/arxiv-2103.01320-CC-BY-4.0.txt. Characters: every retained excerpt is pure ASCII, so the delivered engine, XeLaTeX with its default Latin Modern text font, reports no missing character.
\begin{lem} \label{lem:first_sec_mom_for_ret}
For every quadruple $(A,A',B,B')$ of Borel subsets of $\R_+$, we have that
\begin{equation}
\E\left[\frac{1}{4n^2}\sum_{\ell,j=1}^{2n-1} \mathds{1}_{A \times B}\left(\bm \xi_\ell,\bm s_\ell\right) \mathds{1}_{A' \times B'}\left(\bm \xi_j,\bm s_j\right)\middle| (\bm s_i)_i\right]\xrightarrow[n\to\infty]{a.s.} \E\left[\mathds{1}_{A \times B}(\bm V,\bm U)\right]\E\left[\mathds{1}_{A' \times B'}(\bm V,\bm U)\right]\label{eq:second_mom_ind_ret}.
\end{equation}
\end{lem}

\begin{lem}\label{lem:ind_fnct}
For all $C,C'\in\mathcal{B}(\R_+^2)$, we have that 
\begin{equation}
\E\left[\frac{1}{4n^2}\sum_{\ell,j=1}^{2n-1}\mathds{1}_{C}\left(\bm \xi_\ell,\bm s_\ell\right)\mathds{1}_{C'}\left(\bm \xi_j,\bm s_j\right)\middle| ( \bm s_i)_i\right]\xrightarrow[n\to\infty]{a.s.} \E\left[\mathds{1}_{C}(\bm V,\bm U)\right]\E\left[\mathds{1}_{C'}(\bm V,\bm U)\right]\label{eq:second_mom_ind_bor}.
\end{equation}
\end{lem}

\begin{proof}
We first fix a rectangle $A\times B\in\mathcal{R}(\R_+^2)$ and we consider the set
\begin{equation}
\mathcal{A}_{A\times B}\coloneqq \left\{C\in\mathcal{B}(\R_+^2)\middle |\text{Eq. }\eqref{eq:second_mom_ind_bor} \text{ holds with } C'=A\times B \right\}.
\end{equation}
By \cref{lem:first_sec_mom_for_ret}, we have $\mathcal{R}(\R_+^2)\subseteq\mathcal{A}_{A\times B}\subseteq \mathcal{B}(\R_+^2)$.  If we show that $\mathcal{A}_{A\times B}$ is a monotone class, then we can conclude that $\mathcal{A}_{A\times B}=\mathcal{B}(\R_+^2)$. Indeed, by the monotone class theorem (note that $\mathcal{R}(\R_+^2)$ is closed under finite intersections), we have that
\begin{equation}\label{eq:class_inclusions}
\mathcal{B}(\R_+^2)= \sigma\left(\mathcal{R}(\R_+^2)\right)=\lambda\left(\mathcal{R}(\R_+^2)\right) \subseteq\mathcal{A}_{A\times B}.
\end{equation}
The equality $\mathcal{A}_{A\times B}=\mathcal{B}(\R_+^2)$ implies that \cref{eq:second_mom_ind_bor} holds for every set in $\mathcal{B}(\R_+^2)\times \mathcal{R}(\R_+^2)$. Finally, if we also show that for any fixed Borel set $C^*\in\mathcal{B}(\R_+^2)$, the set
\begin{equation}
\mathcal{A}_{C^*}\coloneqq \left\{C'\in\mathcal{B}(\R_+^2)\middle |\text{Eq. }\eqref{eq:second_mom_ind_bor} \text{ holds with } C=C^* \right\}
\end{equation}
is a monotone class, then using again the same arguments that we used in \cref{eq:class_inclusions} (note that $\mathcal{R}(\R_+^2)\subseteq\mathcal{A}_{C^*}$ thanks to the previous step) we can conclude that \cref{eq:second_mom_ind_bor} holds for every pair of sets in $\mathcal{B}(\R_+^2)\times \mathcal{B}(\R_+^2)$, proving the lemma. Therefore, in order to conclude the proof, it is sufficient to show that $\mathcal{A}_{A \times B}$ and $\mathcal{A}_{C^*}$ are monotone classes.

\medskip


We start by proving that $\mathcal{A}_{A\times B}$ is a monotone class:

\begin{itemize}
\item Obviously $\R_+^2\in \mathcal{A}_{A\times B}$.
\item If $C,D\in\mathcal{A}_{A\times B}$ and $C\subseteq D$, then $D\setminus C\in \mathcal{A}_{A\times B}$ because $\mathds{1}_{D\setminus C}=\mathds{1}_{D}-\mathds{1}_{C}$. 
\item Let now $(C_m)_{m\in\N}$ be a sequence of sets in $\mathcal{A}_{A\times B}$ such that $C_m \subseteq C_{m+1}$ for all $m\in\N$. We want to show that $C\coloneqq\bigcup_m C_m\in \mathcal{A}_{A\times B}$, i.e.\ that 
\begin{equation}\label{eq:ifbueiwbfowebnfoewinf}
\E\left[\frac{1}{4n^2}\sum_{\ell,j=1}^{2n-1}\mathds{1}_{C}\left(\bm \xi_\ell,\bm s_\ell\right)\mathds{1}_{A\times B}\left(\bm \xi_j,\bm s_j\right)\middle| ( \bm s_i)_i\right]
\xrightarrow[n\to\infty]{a.s.} 
\E\left[\mathds{1}_{C}(\bm V,\bm U)\right]\E\left[\mathds{1}_{A \times B}(\bm V,\bm U)\right].
\end{equation}
Since $\mathds{1}_{C}=\lim_{m\to \infty}\mathds{1}_{C_m}$, then by monotone convergence we have for all $n\in \N$,
\begin{multline}\label{eq:lim_of_indicator}
\E\left[\frac{1}{4n^2}\sum_{\ell,j=1}^{2n-1}\mathds{1}_{C_m}\left(\bm \xi_\ell,\bm s_\ell\right)\mathds{1}_{A\times B}\left(\bm \xi_j,\bm s_j\right)\middle| ( \bm s_i)_i\right]
\xrightarrow[m\to\infty]{a.s.}\\
\frac{1}{4n^2}\sum_{\ell,j=1}^{2n-1}\E\left[\mathds{1}_{C}\left(\bm \xi_\ell,\bm s_\ell\right)\mathds{1}_{A\times B}\left(\bm \xi_j,\bm s_j\right)\middle|  \bm s_\ell,\bm s_j\right].
\end{multline}
We also claim that:
\begin{itemize}
	\item[(a)] For all $m\in\N$, 
	\begin{equation}
	\frac{1}{4n^2}\sum_{\ell,j=1}^{2n-1}\E\left[\mathds{1}_{C_m}\left(\bm \xi_\ell,\bm s_\ell\right)\mathds{1}_{A\times B}\left(\bm \xi_j,\bm s_j\right)\middle|  \bm s_\ell,\bm s_j\right]\xrightarrow[n\to\infty]{a.s.} \E\left[\mathds{1}_{C_m}(\bm V,\bm U)\right]\E\left[\mathds{1}_{A \times B}(\bm V,\bm U)\right].
	\end{equation}
	\item[(b)] The convergence in \cref{eq:lim_of_indicator} holds uniformly for all $n\in\N$.
\end{itemize}
Item (a) holds since $C_m \in \mathcal{A}_{A\times B}$.
Item (b) will be proved at the end. Items (a) and (b) allow us to exchange the following a.s.-limits as follows 
\begin{align}
\lim_{n\to \infty}\E&\left[\frac{1}{4n^2}\sum_{\ell,j=1}^{2n-1}\mathds{1}_{C}\left(\bm \xi_\ell,\bm s_\ell\right)\mathds{1}_{A\times B}\left(\bm \xi_j,\bm s_j\right)\middle| ( \bm s_i)_i\right]\\
\stackrel{\eqref{eq:lim_of_indicator}}{=} &\lim_{n\to \infty}\lim_{m\to \infty}
\frac{1}{4n^2}\sum_{\ell,j=1}^{2n-1}\E\left[\mathds{1}_{C_m}\left(\bm \xi_\ell,\bm s_\ell\right)\mathds{1}_{A\times B}\left(\bm \xi_j,\bm s_j\right)\middle|  \bm s_\ell,\bm s_j\right]\\
= &\lim_{m\to \infty}\lim_{n\to \infty}
\frac{1}{4n^2}\sum_{\ell,j=1}^{2n-1}\E\left[\mathds{1}_{C_m}\left(\bm \xi_\ell,\bm s_\ell\right)\mathds{1}_{A\times B}\left(\bm \xi_j,\bm s_j\right)\middle|  \bm s_\ell,\bm s_j\right]\\
=
&\lim_{m\to \infty}\E\left[\mathds{1}_{C_m}(\bm V,\bm U)\right]\E\left[\mathds{1}_{A \times B}(\bm V,\bm U)\right]=\E\left[\mathds{1}_{C}(\bm V,\bm U)\right]\E\left[\mathds{1}_{A \times B}(\bm V,\bm U)\right],
\end{align}
where the last equality follows by monotone convergence. This proves \cref{eq:ifbueiwbfowebnfoewinf} and concludes the proof (up to proving item (b)) that $\mathcal{A}_{A\times B}$ is a monotone class.

\textbf{Proof of  item (b). }Since $\mathds{1}_{C\setminus C_m}=\mathds{1}_{C}-\mathds{1}_{C_m}$, it is enough to show that
\begin{equation}
\sup_{n}\E\left[\frac{1}{4n^2}\sum_{\ell,j=1}^{2n-1}\mathds{1}_{C\setminus C_m}\left(\bm \xi_\ell,\bm s_\ell\right)\mathds{1}_{A\times B}\left(\bm \xi_j,\bm s_j\right)\middle| ( \bm s_i)_i\right]\xrightarrow[m\to\infty]{a.s.}0.
\end{equation}
We set $D_m\coloneqq C\setminus C_m$, and define 
$$(D_m)^{s}\coloneqq\{x\in \R_+|(x,s)\in D_m\},\quad \text{for all} \quad s\in\R,$$
 
$$\pi_Y(D_m)\coloneqq\{y\in \R_+| \exists x\in \R_+ \text{ s.t. }(x,y)\in D_m\}.$$ 

Since  $\bm \xi_\ell\stackrel{d}{=}\bm V$ for all $\ell\in[2n-1]$, then for all $\ell,j\in[2n-1]$, a.s.
\begin{align}
\E\left[\mathds{1}_{D_m}\left(\bm \xi_\ell,\bm s_\ell\right)\mathds{1}_{A\times B}\left(\bm \xi_j,\bm s_j\right)\middle| \bm s_\ell,\bm s_j\right]
&=\P\left(\bm \xi_\ell\in (D_m)^{\bm s_\ell},\bm \xi_j\in A\middle| \bm s_\ell\right)\mathds{1}_{\pi_Y(D_m)}(\bm s_\ell)\mathds{1}_{B}(\bm s_j)\\
&\leq \P\left(\bm \xi_\ell\in (D_m)^{\bm s_\ell}\middle| \bm s_\ell\right)=\P\left(\bm V\in (D_m)^{\bm s_\ell}\middle| \bm s_\ell\right).\label{eq:bound_for_expect1}
\end{align}
Therefore a.s.
\begin{equation}\label{eq:bnd_for_exp}
\sup_{n}\E\left[\frac{1}{4n^2}\sum_{\ell,j=1}^{2n-1}\mathds{1}_{C\setminus C_m}\left(\bm \xi_\ell,\bm s_\ell\right)\mathds{1}_{A\times B}\left(\bm \xi_j,\bm s_j\right)\middle| ( \bm s_i)_i\right]
\leq\P\left(\bm V\in (C\setminus C_m)^{\bm s_\ell}\middle| \bm s_\ell\right).
\end{equation}
Now the sequence $\P\left(\bm V\in (C\setminus C_m)^{\bm s_\ell}\middle| \bm s_\ell\right)$ is a.s.\ non-increasing (because  $C_m \subseteq C_{m+1}$) and hence has an a.s.\ limit. The limit is non-negative and its expectation is the limit of expectations which is $0$ because  $C\coloneqq\bigcup_m C_m$. This completes the proof of item (b).
\end{itemize}


It remains to prove that $\mathcal{A}_{C*}$ is also a monotone class.
The proof is similar to the proof above, replacing the bound in \cref{eq:bound_for_expect1} by
\begin{multline}\label{eq:bound_for_expect2}
\E\left[\mathds{1}_{C^*}\left(\bm \xi_\ell,\bm s_\ell\right)\mathds{1}_{D_m}\left(\bm \xi_j,\bm s_j\right)\middle| \bm s_\ell,\bm s_j\right]\\
=\P\left(\bm \xi_\ell\in (C^*)^{\bm s_\ell},\bm \xi_j\in (D_m)^{\bm s_j}\middle| \bm s_\ell,\bm s_j\right)\mathds{1}_{\pi_Y(C^*)}(\bm s_\ell)\mathds{1}_{\pi_Y(D_m)}(\bm s_j)
\leq \P\left(\bm \xi_j\in (D_m)^{\bm s_j}\middle|\bm s_j\right).
\end{multline}
This completes the proof of the lemma.
\end{proof}

\end{document}
